\section{Localization on curves}
\label{sec:compact-growing-jets}

We now transfer the transition of \cref{prop:degenerate-kernel-transition} to a compact curve. The selected sections concentrate on a shrinking disk around the distinguished point. When the remainder of the metric weight is negligible on that disk, the norm there can be compared directly with the model norm, and restriction to the disk gives an upper bound for the kernel. For the lower bound, positivity of the curvature away from the distinguished point permits a global $L^2$ correction of localized model sections that preserves the growing vanishing order.

\Needspace{15\baselineskip}
Let $C$ be a connected compact Riemann surface, let $x\in C$, and let $(L,h)$ be a holomorphic line bundle on $C$ with a smooth Hermitian metric. Fix an integer $r\geq1$. Suppose that in a holomorphic coordinate $z$ centered at $x$ and a holomorphic frame $e$,
\[
 |e|_h^2=\mathrm{e}^{-\phi},\qquad
 \phi(z)=|z|^{2r}+O(|z|^{2r+1}).
\]
Let $\nu$ be a smooth positive volume form on $C$, written as $a(z)\,\mathrm{d}A$ in this coordinate, where $\mathrm{d}A$ is Euclidean area and $a(0)=1$. For positive integers $m$, choose nonnegative integers $q_m$ such that
\[
 q_m\longrightarrow\infty,\qquad q_m=o(m^{1/(2r+1)}).
\]
Write $\mathfrak m_x$ for the ideal sheaf of $x$, and equip
\[
 V_m=H^0(C,L^m\otimes\mathfrak m_x^{q_m})
\]
with the norm $\|s\|_m^2=\int_C|s|_{h^m}^2\,\mathrm{d}\nu$. Its density is $B_m=\sum_j|s_{m,j}|_{h^m}^2$ for any orthonormal basis $(s_{m,j})$, and $B_m=0$ if $V_m=0$. As in \cref{sec:degenerate-selected-kernels}, $\Phi$ denotes the standard Gaussian distribution function.

\Needspace{16\baselineskip}
\begin{theorem}\label{thm:compact-growing-jets}\index{Bergman kernel!growing jets}
For a fixed real $\tau$, let $x_m$ lie in the coordinate neighborhood and satisfy
\[
 T_m:=m|z(x_m)|^{2r}=q_m/r+\tau\sqrt{q_m/r}.
\]
Then
\begin{equation}\label{eq:compact-growing-jets-upper}
 \limsup_{m\to\infty}
 \frac{\pi B_m(x_m)}{mr^2|z(x_m)|^{2r-2}}\leq\Phi(\tau).
\end{equation}
If the curvature of $(L,h)$ is semipositive on $C$ and strictly positive on $C\setminus\{x\}$, this ratio converges to $\Phi(\tau)$.
\end{theorem}
\begin{proof}
We compare $B_m$ with the model kernels of \cref{sec:degenerate-selected-kernels} on a shrinking disk, which gives the upper bound. For the lower bound, we cut off model polynomials concentrated near $x$ and correct them to global sections by an $L^2$ estimate. Put $z_m=z(x_m)$, $\rho_m=|z_m|$, $R_m=(q_m/(rm))^{1/(2r)}$ and
\[
 S_m=\frac{mr^2\rho_m^{2r-2}}\pi.
\]
Then $R_m\to0$ and $\rho_m/R_m\to1$, and $T_m>0$ for all sufficiently large $m$. Positive constants denoted by $M$, with subscripts when needed, may change between estimates. They are independent of $m$ and of the argument of $z_m$; they may depend on the fixed geometric data, $r$, $\tau$, and fixed auxiliary cutoffs or radii. We write $q=q_m$ when no confusion is possible. The model densities $B_{m,r}^{(q)}$ are those of \cref{sec:degenerate-selected-kernels}.

\Needspace{10\baselineskip}
\emph{Step 1: the kernel on a shrinking disk.}
Fix a real $A>1$ and set $D_m=\{|z|<A R_m\}$. On $D_m$,
\[
 m|\phi(z)-|z|^{2r}|\leq M_A mR_m^{2r+1}
 =M_A\frac{q_mR_m}{r}=o(1),\qquad
 \sup_{D_m}|a-1|=o(1).
\]
The first limit follows from $q_mR_m=r^{-1/(2r)}(q_m^{2r+1}/m)^{1/(2r)}\to0$. Hence, on $D_m$, the norm defined by $h^m$ and $\nu$ and the radial norm with weight $\mathrm{e}^{-m|z|^{2r}}\,\mathrm{d}A$ differ by a factor $1+o(1)$, uniformly over all holomorphic functions on $D_m$. The same comparison holds for the corresponding pointwise weights at $z_m$.

For integers $k\geq0$, let $I_{m,k}$ and $J_{m,k}$ be the squared norms of $z^k$ on $D_m$ and on the whole plane, respectively, with this radial weight. Polar integration gives
\[
 \begin{aligned}
 I_{m,k}&=\frac\pi r m^{-(k+1)/r}\gamma((k+1)/r,U_m),
 &U_m&=m(A R_m)^{2r},\\
 J_{m,k}&=\frac\pi r m^{-(k+1)/r}\Gamma((k+1)/r),
 &\gamma(b,U)&=\int_0^U t^{b-1}\mathrm{e}^{-t}\,\mathrm{d}t.
 \end{aligned}
\]
Here $b,U>0$. By Parseval's identity on circles and Tonelli's theorem, the normalized monomials form a complete orthonormal system in the space of holomorphic functions on $D_m$ with finite radial norm: the squared norm of a Taylor series is the sum of the squared norms of its monomials, so its partial sums converge in that norm. Therefore the density of the subspace of functions vanishing to order at least $q$ is
\[
 B_m^{D}(z_m)=\mathrm{e}^{-T_m}\sum_{k\geq q}\frac{\rho_m^{2k}}{I_{m,k}}.
\]
Point evaluation is continuous by the local submean inequality, so this sum is also the supremum of $\mathrm{e}^{-T_m}|f(z_m)|^2$ over the unit ball of that subspace.

\Needspace{8\baselineskip}
Choose fixed real numbers $1<c<d<A^{2r}$ and put $k_m=\lfloor rcT_m\rfloor$. For large $m$, $k_m\geq q$ and $U_m\geq dT_m$. Uniformly for $0\leq k\leq k_m+1$,
\begin{equation}\label{eq:truncated-gamma-comparison}
 \frac{I_{m,k}}{J_{m,k}}=1+O(\mathrm{e}^{-\eta T_m})
\end{equation}
for some fixed $\eta>0$. Indeed, let $G_b$ be a random variable with Gamma density $t^{b-1}\mathrm{e}^{-t}/\Gamma(b)$ on $(0,\infty)$. A change of variable gives $\mathbb E\mathrm{e}^{sG_b}=(1-s)^{-b}$ for $0<s<1$. With $b=(k+1)/r\leq cT_m+O_r(1)$, Markov's inequality yields
\[
 1-\frac{\gamma(b,U_m)}{\Gamma(b)}
 \leq\mathrm{e}^{-sU_m}(1-s)^{-b}.
\]
Since $d>c$, choose $s>0$ small enough that $sd+c\log(1-s)>0$. The right side is at most $M\mathrm{e}^{-\eta T_m}$, as required.

The Cauchy--Schwarz inequality for the positive radial measure gives
$I_{m,k+1}^2\leq I_{m,k}I_{m,k+2}$, so $I_{m,k+1}/I_{m,k}$ increases with $k$. By \eqref{eq:truncated-gamma-comparison} and the asymptotics of Gamma ratios,
\[
 \rho_m^2\frac{I_{m,k_m}}{I_{m,k_m+1}}
 =(1+o(1))T_m^{1/r}
   \frac{\Gamma((k_m+1)/r)}{\Gamma((k_m+2)/r)}
 \longrightarrow c^{-1/r}<1.
\]
Choose $\kappa\in(c^{-1/r},1)$. For all large $m$, the ratio of consecutive terms $\rho_m^{2k}/I_{m,k}$ is at most $\kappa$ for every $k\geq k_m$. Consequently,
\[
 \mathrm{e}^{-T_m}\sum_{k>k_m}\frac{\rho_m^{2k}}{I_{m,k}}
 \leq\frac{\kappa}{1-\kappa}
          \mathrm{e}^{-T_m}\frac{\rho_m^{2k_m}}{I_{m,k_m}}
 =o(S_m).
\]
For the last equality, replace $I_{m,k_m}$ by $J_{m,k_m}$ using \eqref{eq:truncated-gamma-comparison}. The resulting term is one of the terms of the series for $B_{m,r}^{(k_m)}(z_m)$. In \eqref{eq:degenerate-poisson-comparison} for this kernel, each residue threshold divided by $T_m$ tends to $c>1$, so Chebyshev's inequality makes all the Poisson tails tend to zero, and the comparison error is $O(T_m^{-1/2})$. Hence $B_{m,r}^{(k_m)}(z_m)=o(S_m)$.

For the terms with $q\leq k\leq k_m$, we use \eqref{eq:truncated-gamma-comparison}. Since $I_{m,k}\leq J_{m,k}$ for every $k$, we obtain
\[
 B_{m,r}^{(q)}(z_m)\leq B_m^{D}(z_m)
 \leq(1+o(1))B_{m,r}^{(q)}(z_m)+o(S_m).
\]
By \cref{prop:degenerate-kernel-transition},
\begin{equation}\label{eq:shrinking-disc-transition}
 B_m^{D}(z_m)/S_m\longrightarrow\Phi(\tau).
\end{equation}
The global norm of a selected section dominates its norm on $D_m$. Restriction to $D_m$ and the preceding comparisons of norms therefore give
$B_m(x_m)\leq(1+o(1))B_m^D(z_m)$, also when $V_m=0$. This proves \eqref{eq:compact-growing-jets-upper}.

\Needspace{13\baselineskip}
\emph{Step 2: polynomials with almost all their norm near the point.}
From now on we assume the curvature hypotheses of the last assertion. Fix a real $c>1$ and put $K_m=\lfloor cq_m\rfloor$. Define
\[
 C_m=\sum_{k=q_m}^{K_m}\frac{\rho_m^{2k}}{J_{m,k}},\qquad
 p_m(z)=C_m^{-1/2}\sum_{k=q_m}^{K_m}
                   \frac{\bar z_m^{\,k}z^k}{J_{m,k}}.
\]
\Needspace{6\baselineskip}
For large $m$ these are well defined, $\|p_m\|_{m,r}=1$, and
\begin{equation}\label{eq:model-polynomial-value}
 \mathrm{e}^{-T_m}|p_m(z_m)|^2
 =B_{m,r}^{(q_m)}(z_m)-B_{m,r}^{(K_m+1)}(z_m)
 =(\Phi(\tau)+o(1))S_m.
\end{equation}
Here the first equality holds because $p_m$ is the normalized evaluation vector at $z_m$ of the span of $z^{q_m},\ldots,z^{K_m}$. The omitted kernel $B_{m,r}^{(K_m+1)}(z_m)$ is $o(S_m)$ by \eqref{eq:degenerate-poisson-comparison}, since its Poisson thresholds divided by $T_m$ tend to $c>1$; the second equality then follows from \cref{prop:degenerate-kernel-transition}.

Choose $\delta\in(0,1/2)$ so that the coordinate disk $|z|<\delta$ satisfies
$\phi(z)\geq\beta|z|^{2r}$ for a fixed $\beta\in(1/2,1)$, and such that $a$ is bounded above there. For a sufficiently large fixed $A>1$,
\begin{equation}\label{eq:model-polynomial-localization}
 \int_{A R_m<|z|<\delta}|p_m|^2\mathrm{e}^{-m\phi}a\,\mathrm{d}A=o(1).
\end{equation}
To see this, write $p_m(z)=\sum_{k=q_m}^{K_m}b_{m,k}z^k$, so that
$\sum_k|b_{m,k}|^2J_{m,k}=1$. The lower bound for $\phi$, the orthogonality of monomials on circles and polar integration bound the left side by a fixed constant times
\[
 \sum_{k=q_m}^{K_m}|b_{m,k}|^2J_{m,k}\,
 \beta^{-(k+1)/r}
 \mathbb P\{G_{(k+1)/r}\geq\beta A^{2r}q_m/r\}.
\]
\Needspace{7\baselineskip}
The exponential moment of $G_b$ at $s=1/2$ gives, for $b,U>0$,
\[
 \beta^{-b}\mathbb P\{G_b\geq U\}
 \leq\exp\bigl(-U/2+b\log(2/\beta)\bigr).
\]
Choose $A$ so large that $\beta A^{2r}/2>c\log(2/\beta)+1$. Since $(k+1)/r\leq(cq_m+1)/r$, the last factor in the sum is $O(\mathrm{e}^{-q_m/r})$, uniformly in $k$. This proves \eqref{eq:model-polynomial-localization}. On $|z|<A R_m$, the norm defined by $h^m$ and $\nu$ and the radial model norm differ by a factor $1+o(1)$, as shown in Step~1, and the model norm of $p_m$ on this disk is at most $1$.

Fix a smooth cutoff function $\chi$ with $0\leq\chi\leq1$, compactly supported in $|z|<\delta$ and equal to one near zero, and extend $\chi p_m e^{\otimes m}$ by zero to $C$. The preceding estimates give
\begin{equation}\label{eq:cutoff-polynomial-norm}
 \|\chi p_m e^{\otimes m}\|_m^2\leq1+o(1).
\end{equation}
The numbers $\delta,\beta,A$ and the cutoff $\chi$ are fixed before $m$ tends to infinity.

\emph{Step 3: a global correction with the same vanishing condition.}
Let $f_m=\bar\partial(\chi p_m e^{\otimes m})$. Its support is contained in a fixed annulus on which $\phi$ has a positive lower bound. The Cauchy--Schwarz inequality and the identity $\sum_k|b_{m,k}|^2J_{m,k}=1$ give, for $|z|<\delta<1$,
\[
 |p_m(z)|^2\leq\sum_{k=q_m}^{K_m}\frac{|z|^{2k}}{J_{m,k}}
 \leq M(K_m+1)m^{(K_m+1)/r}.
\]
Here we used that $\Gamma(b)$ has a positive lower bound for $b\geq1/r$. Since $q_m\log m=o(m)$, the factor $\mathrm{e}^{-m\phi}$ on the annulus dominates this growth of the coefficients.

\Needspace{12\baselineskip}
Regard $\nu$ as a K\"ahler form on the curve; it defines pointwise norms of forms and the induced metric $h_{K_C^{-1}}$ on the anticanonical bundle $K_C^{-1}$. Choose a smooth cutoff function $\eta$ with $0\leq\eta\leq1$, supported in the coordinate chart and equal to one near $x$; shrink its support if necessary. For a sufficiently small fixed $t_0>0$, the global smooth functions
\[
 \ell_t=\eta\log(|z|^2+t^2),\qquad 0<t\leq t_0,
\]
extended by zero outside the chart, are nonpositive. Their complex Hessians are nonnegative where $\eta=1$, and elsewhere they are bounded below independently of $t$, because $\eta=1$ near $x$.

Choose a fixed smooth real function $\psi$ that equals a sufficiently large positive multiple of $|z|^2$ near $x$ and is cut off inside the chart, so that
\[
 \mathrm{i}\Theta(K_C^{-1})+\mathrm{i}\partial\bar\partial\psi\geq2\nu
\]
on a neighborhood of $x$ where $\eta=1$. Here $\mathrm{i}\Theta$ denotes the Chern curvature form, without the factor $1/(2\pi)$, so that the weight $\mathrm{e}^{-\psi}$ contributes $\mathrm{i}\partial\bar\partial\psi$. For all large $m$ and all $0<t\leq t_0$, equip
$E_m=L^m\otimes K_C^{-1}$ with the smooth metric
\[
 H_{m,t}=(h^m\otimes h_{K_C^{-1}})\,
                     \mathrm{e}^{-\psi-q_m\ell_t}.
\]
This metric satisfies
\begin{equation}\label{eq:jet-correction-curvature}
 \mathrm{i}\Theta(E_m,H_{m,t})\geq\nu.
\end{equation}
Indeed, near $x$ the two terms involving $h$ and $\ell_t$ are nonnegative, and the other two together are at least $2\nu$. On the complement of a fixed such neighborhood, compactness and strict positivity give a lower bound $c_0m\nu$ for $m\mathrm{i}\Theta(L,h)$, with $c_0>0$. Since $q_m=o(m)$, this term absorbs all remaining negative terms, which are bounded by $M(q_m+1)\nu$.

On the support of $f_m$, the functions $\ell_t$ are bounded below uniformly in $t$. The polynomial estimate, the bounded derivatives of $\chi$, and $q_m\log m=o(m)$ therefore give a constant $a_0>0$ such that, for all large $m$,
\[
 \int_C|f_m|_{h^m,\nu}^2
        \mathrm{e}^{-\psi-q_m\ell_t}\,\mathrm{d}\nu
 \leq\mathrm{e}^{-a_0m}\qquad(0<t\leq t_0).
\]
On a curve, $K_C=\Lambda^{1,0}$, so an $L^m$-valued $(0,1)$-form is the same as an $E_m$-valued $(1,1)$-form, and the induced metrics give it the same pointwise norm; the anticanonical factor in $E_m$, whose curvature enters \eqref{eq:jet-correction-curvature}, accounts for this change of bidegree. In bidegree $(1,1)$ on a curve, \eqref{eq:jet-correction-curvature} makes the curvature operator at least the identity. Since the curve is compact and every $(1,1)$-form on it is $\bar\partial$-closed, the $L^2$ existence theorem for complete K\"ahler manifolds \cite[Theorem~5.1]{DemAM11} applies. It gives an $L^m$-valued weak solution $u_{m,t}$ with
\[
 \bar\partial u_{m,t}=f_m,\qquad
 \int_C|u_{m,t}|_{h^m}^2
       \mathrm{e}^{-\psi-q_m\ell_t}\,\mathrm{d}\nu
 \leq\mathrm{e}^{-a_0m}.
\]

Now fix $m$ and let $t\downarrow0$. Since $\ell_t\leq0$ and $\psi$ is bounded, the solutions $u_{m,t}$ are bounded in the original $L^2$ norm, and a subsequence converges weakly to some $u_m$. Continuity of $\bar\partial$ in the sense of distributions gives $\bar\partial u_m=f_m$. For fixed $t'>0$ and $t<t'$, we have $\ell_t\leq\ell_{t'}$, so the same bound holds in the smooth weighted space with weight $\mathrm{e}^{-\psi-q_m\ell_{t'}}$, whose norm is equivalent to the original one for this fixed $m$ and $t'$. Weak lower semicontinuity, followed by monotone convergence as $t'\downarrow0$, gives
\[
 \int_C|u_m|_{h^m}^2\mathrm{e}^{-\psi-q_m\ell}\,\mathrm{d}\nu
 \leq\mathrm{e}^{-a_0m},\qquad
 \ell=\eta\log|z|^2.
\]
In particular, $\|u_m\|_m^2\leq M\mathrm{e}^{-a_0m}$. Near $x$ the right side $f_m$ vanishes, so the Cauchy--Riemann equation, in the sense of distributions, makes $u_m$ holomorphic in the frame $e^{\otimes m}$. The smooth factors in the weight are bounded above and below near $x$ for each fixed $m$, so the coefficient of $u_m$ is square integrable against $|z|^{-2q_m}$. By \cref{lem:smooth-center-multiplier-ideal} in dimension one, this forces vanishing order at least $q_m$ at $x$.

Since $p_m$ also vanishes to order at least $q_m$ at zero, it follows that $s_m=\chi p_m e^{\otimes m}-u_m$ is a global holomorphic section in $V_m$. By \eqref{eq:cutoff-polynomial-norm}, $\|s_m\|_m\leq1+o(1)$.

\emph{Step 4: evaluation at the moving point.}
For large $m$, the disk $D_m$ lies in the region where $\chi=1$ and $f_m=0$. Apply the local comparison of norms and the extremal property of $B_m^D$ to the holomorphic function representing $u_m$ on $D_m$, which vanishes to order at least $q_m$. By \eqref{eq:shrinking-disc-transition},
\[
 |u_m(x_m)|_{h^m}^2
 \leq(1+o(1))B_m^D(z_m)\|u_m\|_m^2
 \leq M S_m\mathrm{e}^{-a_0m}.
\]
The pointwise factor $\mathrm{e}^{-m\phi(z_m)}$ differs from $\mathrm{e}^{-T_m}$ by a factor $1+o(1)$. Since $\Phi(\tau)>0$, \eqref{eq:model-polynomial-value} and the triangle inequality give
\[
 |s_m(x_m)|_{h^m}^2\geq(\Phi(\tau)+o(1))S_m.
\]
In particular, these sections are nonzero for large $m$. The extremal characterization of $B_m$ and $\|s_m\|_m\leq1+o(1)$ yield the matching lower bound, which proves the limit.
\end{proof}

The hypothesis $q_m=o(m^{1/(2r+1)})$ enters only through $mR_m^{2r+1}=o(1)$, which makes the first remainder term of the metric weight negligible throughout the disk carrying the norm. For larger orders, the curvature quadrature domains of \cref{sec:curvature-quadrature-jets} absorb the harmonic part of this error and determine where the transition occurs. \Cref{thm:curvature-boundary-jets} gives the corrected transition for all sublinear orders, and \cref{thm:optimal-growing-jet-range} determines the optimal range in which the uncorrected law holds for every metric. On a fibration over a curve, the kernel identities of \cref{sec:partial-submersion} transfer these transitions whenever the descended spaces carry the same vanishing conditions and the induced norm.

\subsection{Curvature quadrature domains}
\label{sec:curvature-quadrature-jets}

Retain $C,x,(L,h),r,z,e$ and $\nu=a(z)\,\mathrm{d}A$ from \cref{sec:compact-growing-jets}, and assume that the curvature is semipositive on $C$ and strictly positive on $C\setminus\{x\}$. We now allow integers $q=q_m\to\infty$ with $q=o(m)$. Write $B_{m,q}$ for the density of $H^0(C,L^m\otimes\mathfrak m_x^q)$ with the same norm. Put
\[
 \rho=\phi_{z\bar z},\qquad t=q/m,\qquad R=(t/r)^{1/(2r)}.
\]
The disc comparison above requires the weight error to be small throughout the disc. A larger range is obtained by absorbing the harmonic part of the error into the holomorphic section and locating the boundary on which the remaining potential has first-order contact. Its curvature quadrature identity determines that boundary.

\Needspace{18\baselineskip}
\begin{lemma}\label{lem:curvature-quadrature-domain}\index{Quadrature domain!curvature}
Let $Q(R,\zeta)$ be jointly smooth on a fixed disk of radius $b>1$,
with $Q(0,\zeta)=|\zeta|^{2r}$. Suppose every $Q_R$ is subharmonic
and $Q_{R,\zeta\bar\zeta}>0$ in a fixed collar of the unit circle.
For small $R$, the largest subharmonic minorant of $Q_R$ with a pole
$2r\log|\zeta|+O(1)$ at zero has exactly one noncontact component
$\Omega_R$. Its boundary is a radial graph near the unit circle. For
every fixed finite $k$ the normalized interior map $F_R$ and the
logarithm-subtracted harmonic part depend smoothly on $R$ in the
corresponding H\"older spaces, with $O(R)$ distance from their radial
values. All estimates used below require only finitely many such norms.
\end{lemma}
\begin{proof}
Write $\mathbb D$ for the unit disc. Fix $k\geq3$ and $0<\alpha<1$. Let $A^{k,\alpha}$ be the real Banach
space of functions holomorphic on $\mathbb D$, $C^{k,\alpha}$ on its closure,
whose value at zero is real. For $h\in A^{k,\alpha}$ near zero set $f(u)=u(1+h(u))$. Small $h$ gives a
univalent $f$, with $f(0)=0$ and $f'(0)>0$. On $|u|=1$ define
\[
 K(R,h)=f(u)(Q_R)_\zeta(f(u))-r,
 \qquad T(R,h)=(P_-K,\operatorname{Re}[K]_0),
\]
where $P_-$ projects onto strictly negative Fourier modes and $[\cdot]_0$
is the constant coefficient. The boundedness of Cauchy projections on
these H\"older spaces follows from the circle Hilbert kernel: subtract
the value at the evaluation point, split its integral at twice the
increment length, and use the H\"older difference bound on the inner
part and the differentiated kernel bound on the outer part. This gives
the $C^{0,\alpha}$ bound; commuting angular derivatives gives the
$C^{k,\alpha}$ bound.
The smooth superposition operator has no derivatives
of $h$. Thus $T$ is a smooth map to the negative-mode space times $\mathbb R$.
At $(R,h)=(0,0)$, $T=0$, and
\[
 D_hK(0,0)v=r^2(v+\bar v),\qquad
 D_hT(0,0)v=(r^2\overline{v-v(0)},\,2r^2v(0)).
\]
This is a bounded real-linear isomorphism. The Banach implicit-function
theorem therefore gives a unique normalized branch $f_R$, smoothly
depending on $R$ in the fixed space, with $\|f_R-\mathrm{id}\|=O(R)$.
Branches obtained in higher fixed spaces agree by lower-space uniqueness.
This supplies all finite spatial estimates and all finite parameter
expansions needed in the asymptotic argument. Each prescribed finite
order has its own sufficiently small parameter interval; no common
interval is inferred by intersecting infinitely many intervals.
The equation has not omitted an imaginary residue condition. The
function $K$ now extends holomorphically to $\mathbb D$, and $K(0)$ is purely
imaginary. Stokes' formula gives
\[
 \frac1{2\pi\mathrm{i}}\int_{|u|=1}(Q_R)_\zeta(f_R(u))f_R'(u)\,\mathrm{d}u
 =\frac1\pi\int_{f_R(\mathbb D)}(Q_R)_{\zeta\bar\zeta}\mathrm{d}A\in\mathbb R.
\]
The residue of $(r+K)f_R'/f_R$ is $r+K(0)$, so $K(0)$ is real as well,
and is zero. Consequently
\[
 S_R(\zeta)=\frac{r+K(f_R^{-1}(\zeta))}{\zeta}
 =r/\zeta+H_R(\zeta)
\]
has a holomorphic $H_R$ on $\Omega_R=f_R(\mathbb D)$ and equals $(Q_R)_\zeta$
on its boundary. Let
\[
 U_R=2r\log|\zeta|+2\operatorname{Re}\int_0^\zeta H_R(w)\,\mathrm{d}w+C_R.
\]
Its complex gradient agrees with that of $Q_R$ on the boundary.
The boundary difference is constant; choose $C_R$ to make it zero.
The difference and its gradient vanish there. Its inward normal second
derivative is $\Delta Q_R>0$, because the tangential second derivative
vanishes. Hence $Q_R-U_R>0$ in a uniform inner collar. At $R=0$,
\[
 U_0=1+2r\log|\zeta|,\qquad
 |\zeta|^{2r}-1-2r\log|\zeta|>0\quad(0<|\zeta|<1).
\]
Uniform convergence of the gap preserves its positive minimum on the remaining compact interior
sets; the logarithmic term controls a neighborhood of zero.

Patch $U_R=Q_R$ off $\Omega_R$. Value and gradient matching eliminate
boundary measures in its Laplacian. Its curvature inside is $r\delta_0$
and outside is $\ddc Q_R\geq0$, so it is an admissible minorant. If $v$
is any competitor, $v-U_R$ is subharmonic on $\Omega_R\setminus\{0\}$,
bounded above at zero and nonpositive on its boundary. Removable
subharmonic extension and the maximum principle give $v\leq U_R$.
Outside the hole this follows from $v\leq Q_R$. This identifies the
envelope and excludes extra noncontact components. The radial graph and
the harmonic-part estimates follow from the normalized map and its
parameter dependence. We henceforth denote this map by $F_R=f_R$.
\end{proof}

Taylor's integral remainder gives, for small signed $R$,
\[
 Q_R(\zeta)=R^{-2r}\phi(R\zeta)
 =|\zeta|^{2r}+R E(R,\zeta)
\]
with jointly smooth $E$. This applies the lemma to the actual weight.
Choose the harmonic primitive with $G_R(0)$ real, and write its envelope inside the hole as
\[
 U_R=2r\log|\zeta|+2\operatorname{Re} G_R(\zeta),\qquad
 G_R=1/2+O(R).
\]
Set $D_t=R\Omega_R$. In physical coordinates the obstacle potential is
$\phi_t(z)=R^{2r}U_R(z/R)$ inside $D_t$ and $\phi$ outside.
Replacing $\phi$ by $\phi_t$ in this chart and keeping $h$ elsewhere defines a global singular metric on $L$. Its potential is $C^1$ across the boundary, is bounded above by $\phi$, and has pole $2t\log|z|+O(1)$.
On the boundary $\phi_z=t/z+g_t'$ for a holomorphic $g_t$. Green's formula gives
\begin{equation}\label{eq:curvature-quadrature-identity}
 \int_{D_t}\rho(z)f(z)\mathrm{d}A=\pi t f(0)
\end{equation}
for every $f$ holomorphic near $\overline{D_t}$. The curvature is $t\delta_x$ inside and the original curvature outside. With $c_R$ a positive periodic function, the boundary is written as $\Gamma_t=\{R c_R(\theta)\mathrm{e}^{\mathrm{i}\theta}\}$, where $c_R-1=O(R)$ in each fixed norm used below.

\subsection{Gaussian transitions at curvature boundaries}
\label{sec:mesoscopic-growing-jets}

The singular metric supported at the center enforces the growing vanishing order in the global correction. Its smooth part is flat inside the quadrature domain, so a bounded function at scale $R$ must restore strict positivity before the $L^2$ estimate is applied.

\Needspace{20\baselineskip}
\begin{lemma}\label{lem:shrinking-order-correction}
For all sufficiently large $m$ there is a uniformly bounded smooth function $b_R$ such that the metric on $L^m\otimes K_C^{-1}$ induced by the singular potential $\psi_m=m\phi_t+b_R$ and the inverse canonical metric of $\nu$ has curvature at least $(c/R^2)\nu$ globally. On exterior annuli of radii
comparable to $R$ the lower bound is $(cq/R^2)\nu$. If an $L^m$-valued
$\bar\partial$ source $\alpha$ is supported on such an exterior annulus,
it has a linear solution $v$ with
\begin{equation}\label{eq:order-preserving-correction}
 \|v\|_{h^m}^2\leq C\frac{R^2}{q}\|\alpha\|_{h^m,\nu}^2.
\end{equation}
If the source vanishes near $x$, $v$ is holomorphic there and vanishes
to order at least $q$.
\end{lemma}
\begin{proof}
Take $b_R=|z|^2/R^2$ on a disk containing the hole, and cut it off before
radius $3R$, with its negative curvature supported outside the hole.
The bump is uniformly bounded and its Hessian is $O(R^{-2})$. On that
support the curvature $m\rho$ is comparable to $q/R^2$, so $q\to\infty$
absorbs the negative part. Inside the hole the bump supplies $c/R^2$.
Elsewhere in a fixed coordinate disk the original curvature is bounded
below by a fixed multiple of $m|z|^{2r-2}$; off that disk it is at least
$cm$. Since $mR^2\to\infty$, all these bounds dominate $c/R^2$ and the
bounded Ricci term of $\nu$.

Apply the pointwise inverse-curvature $L^2$ theorem for singular metrics,
in the formulation recorded in \cref{sec:history-sections}.
For sections use $L^m\otimes K_C^{-1}$ in bidegrees $(1,0)$ and $(1,1)$;
this includes the Ricci term and preserves the original form norms under
the canonical identification. More explicitly, if $H$ is the singular metric on $L^m$ and the twisted curvature is at least $\lambda\nu$ for a continuous $\lambda>0$, the estimate is
\[
 \int_C|v|_H^2\,\mathrm{d}\nu
 \leq\int_C\lambda^{-1}|\alpha|_{H,\nu}^2\,\mathrm{d}\nu.
\]
The inverse curvature on the support is
$O(R^2/q)$. There $\phi_t=\phi$, and the bump is bounded. Choosing the minimal solution in the singular weighted Hilbert space makes the correction linear in its source. Elsewhere
$\mathrm{e}^{-m\phi}\leq C \mathrm{e}^{-\psi_m}$. These comparisons yield \eqref{eq:order-preserving-correction}.
To use a continuous positive curvature lower bound, interpolate from
$c/R^2$ inside the hole to $cq/R^2$ across a thin exterior collar.
For the separated cutoff supports below, choose this collar to miss
the source. More generally, for an exterior source accumulating on
$\Gamma_t$, shrink the collar to the boundary. Dominated convergence
in the inverse-curvature integrals and weak compactness in the fixed
weighted Hilbert space give the same estimate; the boundary has area
zero. No discontinuous lower bound is required by the cited theorem.
Approximation of the positive singular metric gives the same estimate
and a weak solution. The distributional Cauchy--Riemann equation makes it holomorphic where the source vanishes.
Near $x$, weighted integrability is against
$|z|^{-2q}$ times a smooth positive factor. A holomorphic term $z^j$ is
integrable precisely when $j\geq q$; thus the solution has the required
vanishing order.
\end{proof}

\Needspace{20\baselineskip}
\begin{theorem}\label{thm:curvature-boundary-jets}\index{Bergman kernel!curvature boundary transition}
For the domains $D_t$ constructed above, let $d_t(z)$ be signed
Euclidean distance from $\Gamma_t$, positive on the exterior, and let
$p_t(z)$ be the nearest boundary point. For every fixed $M>0$, uniformly
in all directions and $|d_t(z)|\leq MR/\sqrt q$,
\begin{equation}\label{eq:curvature-boundary-transition}
 \frac{\pi a(z)B_{m,q}(z)}{m\rho(z)}
 -\Phi\bigl(2\sqrt{m\rho(p_t(z))}\,d_t(z)\bigr)\longrightarrow0.
\end{equation}
If $d_t(z)=o(R)$ and $\sqrt q\,d_t(z)/R$ tends to $-\infty$ or
$+\infty$, the same normalized kernel tends to $0$ or $1$, respectively.
Thus the corrected center and transition length are $\Gamma_t$ and
$1/(2\sqrt{m\rho})$, for the entire sublinear range.
\end{theorem}
\begin{proof}
Define on the unit disk
\[
 \begin{aligned}
 A_R(w)&=G_R(F_R(w))+r\log(F_R(w)/w),\\
 \Psi_R&=Q_R(F_R(w))-2\operatorname{Re} A_R(w)-2r\log|w|.
 \end{aligned}
\]
The logarithm uses the branch near one. Here $A_R=1/2+O(R)$ and
$F_R=w+O(R)$ in every fixed norm used below. The effective potential is
positive inside, with zero value and gradient on the boundary. For
$w=\mathrm{e}^{v+\mathrm{i}\theta}$ its expansion is
\begin{equation}\label{eq:curvature-ring-expansion}
 \Psi_R=(2r^2+O(R))v^2+O(v^3).
\end{equation}
The exact quadratic coefficient is
$2(Q_R)_{\zeta\bar\zeta}(F_R(\mathrm{e}^{\mathrm{i}\theta}))|F_R'(\mathrm{e}^{\mathrm{i}\theta})|^2$.

The maps need not extend holomorphically across the circle. To obtain
genuine exterior charts, truncate their Taylor series at
$K=\lfloor q^{1/16}\rfloor$, denoting them by $F_K,A_K$.
Smooth Fourier coefficient bounds in a sufficiently high fixed norm
make the truncation error $O(q^{-10})$, including the fixed derivatives
needed here. On $|w|\leq1+c/K$ their derivatives remain bounded and
$F_K-w,A_K-1/2=O(R)$: the factor $(1+c/K)^K\leq \mathrm{e}^c$ controls the sums.
In particular $F_K$ is univalent there for small $R$. With
$\delta=c/K$ and
\[
 \Psi_K=Q_R(F_K)-2\operatorname{Re} A_K-2r\log|w|,
\]
value and first derivative errors on the circle are $O(q^{-10})$.
Uniformly for $|v|\leq\delta$,
\begin{equation}\label{eq:polynomial-ring-expansion}
 \Psi_K=(2r^2+O(R)+O(|v|))v^2+O(q^{-10})(1+|v|).
\end{equation}
On the interior the effective potential has uniform positive quadratic
growth near the circle, a positive gap on compact annuli farther inside,
and logarithmic growth at zero.

In $z=RF_K(w)$ a local holomorphic section of order at least $q$ is
uniquely represented, in the original frame, by
\begin{equation}\label{eq:quadrature-section-gauge}
 s=w^q\exp((q/r)A_K(w))f(w),\qquad
 |s|_{h^m}^2=|f(w)|^2\mathrm{e}^{-(q/r)\Psi_K(w)},
\end{equation}
where $f$ is holomorphic at zero and has only nonnegative Taylor powers.
The volume in this chart is $R^2(1+O(R))\mathrm{d}A_w$. Put
$\eta=\sqrt{rq}\,v$. On each fixed $|\eta|\leq L$,
\[
 (q/r)\Psi_K=2\eta^2+o(1).
\]
For an observation $\eta_0$ in a fixed bounded interval, restriction to
this annulus bounds the kernel from above. Its limiting diagonal
monomial norms, for $f(w)=\sum_{j\geq0}b_jw^j$, are
\[
 n_{j,L}=\frac{2\pi R^2}{\sqrt{rq}}I_L(j/\sqrt{rq}),\qquad
 I_L(s)=\int_{-L}^L \mathrm{e}^{2s\eta-2\eta^2}\,\mathrm{d}\eta.
\]
After normalization by $rq/(\pi R^2)$, the limiting upper kernel is
\begin{equation}\label{eq:positive-mode-ring-kernel}
 \frac12\mathrm{e}^{-2\eta_0^2}\int_0^\infty\frac{\mathrm{e}^{2s\eta_0}}{I_L(s)}\,\mathrm{d}s.
\end{equation}
The sum-to-integral passage uses mesh $1/\sqrt{rq}$. To control its tail,
integrate $I_L$ on a fixed exterior subinterval beyond every observation
$\eta_0$; the resulting exponential bound is uniform in the mesh.
Thus Riemann sums and then $L\to\infty$ are justified. Since
$I_\infty(s)=\sqrt{\pi/2}\,\mathrm{e}^{s^2/2}$, \eqref{eq:positive-mode-ring-kernel} tends to
\begin{equation}\label{eq:gaussian-ring-limit}
 \frac1{\sqrt{2\pi}}\int_0^\infty \mathrm{e}^{-(s-2\eta_0)^2/2}\,\mathrm{d}s
 =\Phi(2\eta_0).
\end{equation}

For the lower bound fix $T$ and use \eqref{eq:quadrature-section-gauge} with
$f$ in the span of $w^j$, $0\leq j\leq T\sqrt{rq}$. Let $\|s\|_{\mathrm{trial}}$ denote the Gaussian trial norm below. Cut off with a radial $\chi$ supported through
$1+\delta$, equal to one through $1+\delta/2$.
The cutoff annulus lies outside the true hole. Equations
\eqref{eq:polynomial-ring-expansion} and \eqref{eq:quadrature-section-gauge} bound the squared source norm,
as a quadratic form on the entire span, by
\[
 \|\bar\partial(\chi s)\|_{h^m,\nu}^2\leq C(R\delta)^{-2}\mathrm{e}^{-cq\delta^2}\|s\|_{\mathrm{trial}}^2.
\]
The linear term $2jv$ is absorbed because $\sqrt q\,\delta\to\infty$.
\Cref{lem:shrinking-order-correction} gives a linear correction $v_s$ satisfying
\begin{equation}\label{eq:shrinking-cutoff-correction}
 \|v_s\|_{h^m}^2\leq C(q\delta^2)^{-1}\mathrm{e}^{-cq\delta^2}\|s\|_{\mathrm{trial}}^2.
\end{equation}
The factors depend on $m$ through $R$ and $q$; the factor $R^2$ in the correction estimate cancels the cutoff derivative. This is needed even when $q$ grows slower than every power of $m$.

The trial Gram form converges in relative operator norm to the diagonal
form
\begin{equation}\label{eq:gaussian-ring-gram}
 n_j=\frac{2\pi R^2}{\sqrt{rq}}\sqrt{\frac\pi2}
          \exp(j^2/(2rq)).
\end{equation}
Indeed the weight ratio is uniform on each fixed $\eta$-band. Outside
it, radial domination and angular Parseval bound arbitrary combinations
by diagonal tails $\mathrm{e}^{-c\eta^2+2T|\eta|}$. First pass to the limit on
the fixed band and then enlarge it. Interior compact gaps and the
logarithmic singularity provide the remaining domination. This controls
the whole quadratic form without a dimension factor.

At points of radius comparable to $R$, a local submean estimate gives
\begin{equation}\label{eq:shrinking-point-evaluation}
 |u(z)|_{h^m}^2\leq C(q/R^2)\|u\|_{h^m}^2
\end{equation}
for a section $u$ holomorphic near the point on a disk of radius
$R/\sqrt q$. Remove the harmonic linear and quadratic Taylor terms
there. The scaled remaining quadratic coefficient is bounded and the
third-order error is $O(q^{-1/2})$, which proves the estimate.
It applies to the corrections, since their sources lie on the distant
cutoff annulus. Together with \eqref{eq:shrinking-cutoff-correction}, it shows that the
corrected global trial kernel differs by $o(q/R^2)$ from the local trial
kernel. The normalized lower limit is \eqref{eq:gaussian-ring-limit} with upper
limit $T$ instead of infinity; let $T\to\infty$.

At a nearest boundary point,
\[
 \begin{aligned}
 v&=\frac{d_t(z)}{R|F_R'(\mathrm{e}^{\mathrm{i}\theta})|}
       +O(d_t(z)^2/R^2)+o(q^{-1/2}),\\
 m\rho(p)R^2|F_R'(\mathrm{e}^{\mathrm{i}\theta})|^2&=rq(1+O(R)).
 \end{aligned}
\]
Also $m\rho(z)/a(z)=(rq/R^2)(1+o(1))$ in the fixed transition band.
These identities convert $2\eta_0$ to the argument in \eqref{eq:curvature-boundary-transition}
and give its exact density normalization.

\emph{Full density and saturation.}
At points $|z|\asymp R$ the full, unselected kernel obeys
\begin{equation}\label{eq:mesoscopic-full-density}
 B_m^{\rm full}(z)=\frac{m\rho(z)}{\pi a(z)}(1+o(1)),
\end{equation}
uniformly in an $o(R)$ neighborhood of $\Gamma_t$. Here is a direct
localization argument, which also supplies the saturation part of
\cref{thm:curvature-boundary-jets}. Remove the harmonic Taylor terms of $m\phi$
at the observation and rescale by $(m\rho(z))^{-1/2}$. On every fixed
rescaled disk the remaining weight tends to $|u|^2$ and the volume to
the constant $a(z)$. Taylor's estimate is uniform because the third
derivatives at radius $R$ are $O(R^{2r-3})$, so the normalized third
order error is $O(q^{-1/2})$. Restriction to such disks bounds the
evaluation norm from above; their Gaussian constant-function norm tends
to $\pi$ as their radius grows. This proves the upper bound in
\eqref{eq:mesoscopic-full-density}.

For the lower bound use the holomorphic exponential cancelling the
harmonic Taylor terms, and cut it off on a disk of radius $\varepsilon R$
centered at the observation, where $\varepsilon\to0$ and
$\sqrt q\,\varepsilon\to\infty$. Its local squared norm is
$(\pi a(z)/(m\rho(z)))(1+o(1))$ relative to its unit weighted value
at the center. Quadratic lower domination bounds the source annulus by
an exponential $\mathrm{e}^{-cq\varepsilon^2}$. Add the same bounded bump of
scale $R$ to the original metric and use the pointwise $\bar\partial$
estimate. The relative squared correction is bounded by
$C(q\varepsilon^2)^{-1}\mathrm{e}^{-cq\varepsilon^2}=o(1)$, after decreasing
$c$ if needed. Its value at the center is $o(1)$ by the local submean
estimate \eqref{eq:shrinking-point-evaluation}. This proves the matching lower bound.

For interior saturation, use the polynomial chart above and put
$u=1-|w|>0$. Restrict the norm in \eqref{eq:quadrature-section-gauge} to the exterior
half-band $0\leq v\leq c/\sqrt q$. Every nonnegative Taylor monomial
of $f$ has norm at least $cR^2/\sqrt q$. Angular Parseval and
Cauchy--Schwarz give
\[
 B_{m,q}(z)\leq
 C\frac{\sqrt q}{R^2}\frac{\mathrm{e}^{-cqu^2}}{1-|w|^2}.
\]
If $u\sqrt q\to\infty$, division by $q/R^2$ makes this tend to zero.
The chart error is $O(q^{-10})$ and does not alter the hypothesis or
the conclusion.

For exterior saturation, use a Gaussian peak at $z$ and cut it off on
a disk of radius
\[
 \varepsilon R,\qquad
 \varepsilon=\min\{d_t(z)/(3R),q^{-1/4}\}.
\]
It is wholly outside the hole. Again $\varepsilon\to0$ and
$\sqrt q\varepsilon\to\infty$. Apply \cref{lem:shrinking-order-correction} using the
singular obstacle metric, which equals the original metric on this
source. The correction has negligible norm and value and vanishes to
order $q$ at $x$. Thus a member of $V_{m,q}$ attains the full Gaussian
density asymptotically. Equation~\eqref{eq:mesoscopic-full-density} is its upper bound.
This proves saturation to one.
\end{proof}

\subsection{Growth ranges and compact examples}
\label{sec:optimal-growing-jet-ranges}

The quadrature boundary differs from the model circle by $O(R^2)$ in physical coordinates. Comparing this displacement with the transition width $R/\sqrt q$ gives the universal range. A holomorphic change of coordinate produces a compact example attaining this displacement.

\Needspace{24\baselineskip}
\begin{theorem}\label{thm:optimal-growing-jet-range}\index{Bergman kernel!optimal growing-jet range}
Let $z_m$ denote the coordinate values of the observation points. For each fixed integer $r\geq1$ and real $\tau$, the law
\begin{equation}\label{eq:optimal-growing-jet-law}
 \frac{\pi B_{m,q}(z_m)}{mr^2|z_m|^{2r-2}}
 \longrightarrow\Phi(\tau),\qquad
 m|z_m|^{2r}=q/r+\tau\sqrt{q/r},
\end{equation}
holds for every choice of the preceding geometric data and every sequence of observation
directions when
$q=o(m^{1/(r+1)})$. This is the optimal universal range: for every
excluded sequence $q\to\infty$, $q=o(m)$, a compact example satisfying
the same hypotheses violates the law along a subsequence.


Write $\Gamma_t$ as $|z|=R c_R(\theta)$. For a fixed metric, the exact
criterion for \eqref{eq:optimal-growing-jet-law} for every sequence of directions and every
fixed $\tau$ is
\begin{equation}\label{eq:fixed-weight-jet-criterion}
 \sqrt q\,\|c_R-1\|_\infty\longrightarrow0.
\end{equation}
\end{theorem}
\begin{proof}
By \cref{lem:curvature-quadrature-domain}, $c_R-1=O(R)$ in every fixed norm required
here. The observation in \eqref{eq:optimal-growing-jet-law} satisfies
\[
 |z_m|/R=1+\tau/(2\sqrt{rq})+O(q^{-1}).
\]
Since the boundary slope tends to zero, its corrected normal argument
is $\tau+o(1)$ whenever $\sqrt q R\to0$. This condition is precisely
$q=o(m^{1/(r+1)})$. Moreover $a(z)=1+o(1)$ and
$\rho(z)=r^2|z|^{2r-2}(1+o(1))$ there, so the two normalizations agree.
The edge theorem proves universal sufficiency.

The same reasoning proves the sufficiency of \eqref{eq:fixed-weight-jet-criterion}.
For necessity choose, on a subsequence, an angle realizing the maximum
boundary displacement and observe at $|z|=R$ with $\tau=0$. A nonzero
finite subsequential scaled displacement gives a nonzero Gaussian
argument; an infinite one gives saturation to zero or one. Radial and
normal distances have relative error tending to zero, since the graph
has vanishing slope. In either case the limit differs from $1/2$.

We give a compact model for sharpness without assumptions at infinity.
On $\mathbb P^1$ choose a radial metric on $\mathcal O(d)$ with local potential
$f(u)$, $u=|w|^2$, equal to $u^r$ near zero, and with
$g(u)=uf'(u)$ strictly increasing on $(0,\infty)$ and tending to $d$.
For large $u$ require $g(u)=d-\epsilon/u$, with $\epsilon>0$.
Such a smooth interpolation exists for a sufficiently large integer
$d$: prescribe a positive $g'$ equal to $r^2u^{r-1}$ near zero and
$\epsilon/u^2$ near infinity, with total integral $d$.
Then $f=d\log u+\mathrm{constant}+\epsilon/u$ near infinity; the
infinity-frame potential is smooth with strictly positive curvature.
Choose a smooth radial volume equal to $\mathrm{d}A_w$ near zero.

For sufficiently small $t$, the radial hole lies in the exact-model
region and is $|w|<R$: rotational invariance and uniqueness make it a
disk, and its curvature mass is $rR^{2r}=t$. The corrected transition
and saturation theorem therefore apply to this compact example.

Change the allowed observation coordinate to $w=z+c z^2$, $c>0$.
The same compact metric now has
$\phi(z)=|z|^{2r}+O(|z|^{2r+1})$, and the volume still has $a(0)=1$.
Vanishing order is invariant. At the raw observation $z=R$,
\[
 2\sqrt{mr^2R^{2r-2}}\,(|w|-R)
 =2c\sqrt r\,\sqrt q R.
\]
The two curvature-model denominators have ratio tending to one.
Every excluded sequence has a subsequence on which $\sqrt q R$ tends
to a positive finite number or to infinity. \Cref{thm:curvature-boundary-jets}, using its saturation conclusion when the
argument diverges, then gives a value strictly larger than $1/2$. This proves sharpness.
\end{proof}

\subsection{Corrections from the smooth weight}
\label{sec:growing-jet-corrections}
\index{Bergman kernel!effective Taylor correction}

Write
\[
 \phi(z)=|z|^{2r}+P_{2r+1}(z,\bar z)+O(|z|^{2r+2}),\qquad
 P(\theta)=P_{2r+1}(\mathrm{e}^{\mathrm{i}\theta},\mathrm{e}^{-\mathrm{i}\theta}),
\]
where $P_{2r+1}$ is a real homogeneous polynomial. Let $\Lambda$ be the
circle Dirichlet-to-Neumann operator,
$\Lambda \mathrm{e}^{\mathrm{i} n\theta}=|n|\mathrm{e}^{\mathrm{i} n\theta}$. The first boundary coefficient is
\begin{equation}\label{eq:first-boundary-coefficient}
 c_R=1+Rb_1+O(R^2),\qquad
 b_1=-\frac{((2r+1)-\Lambda)P}{4r^2}.
\end{equation}
To derive it, write the interior obstacle to first order as
$1+2r\log s+R H$, with harmonic $H$. Value matching on
$s=1+Rb_1$ gives $H|_{s=1}=P$; the terms involving $b_1$ cancel since
the radial reference already matches first derivatives. Derivative
matching gives $4r^2b_1+(2r+1)P-\Lambda P=0$.
Pure harmonic terms disappear, as required by invariance under frame
changes.

For $q=O(m^{1/(r+1)})$, at the raw observation the corrected argument is
\begin{equation}\label{eq:first-growing-jet-shift}
 \tau+\frac{\sqrt q R}{2r^{3/2}}((2r+1)-\Lambda)P(\theta)+o(1),
\end{equation}
uniformly in $\theta$. The kernel differs from $\Phi$ of the displayed
argument by $o(1)$. The error is small because $\sqrt q R^2\to0$ in
this range. For $w=z+cz^2$, $P=2r\operatorname{Re}(c\mathrm{e}^{\mathrm{i}\theta})$ and the shift is
$2\sqrt{rq}R\operatorname{Re}(c\mathrm{e}^{\mathrm{i}\theta})$, agreeing with the compact example.

For every fixed integer $J\geq1$ the parameter construction gives
\[
 c_R=1+\sum_{j=1}^J R^j b_j+O(R^{J+1})
\]
in every fixed norm needed. There is an explicit recursion: expand the
interior harmonic obstacle in powers of $R$. At order $j$, value
matching determines the boundary data of its new harmonic term from
the degree $2r+j$ Taylor term and all preceding coefficients. Harmonic
extension determines this term. Derivative matching is then
$4r^2b_j+\text{known forcing}=0$. The new $b_j$ does not enter the
order-$j$ value equation. Nonlinear lower-order terms contribute when
$j\geq2$.

The truncated curve changes the physical center by $O(R^{J+2})$.
Using it in the geometric correction is valid whenever
\[
 \sqrt q R^{J+1}\to0,
 \quad\hbox{equivalently}\quad
 q=o(m^{(J+1)/(r+J+1)}).
\]
If $b_1=\cdots=b_{j-1}=0$ but $b_j\not\equiv0$, the exact criterion
\eqref{eq:fixed-weight-jet-criterion} makes the fixed-weight uncorrected range
$q=o(m^{j/(r+j)})$, sharply in the all-direction sense. Higher terms
are relatively $O(R)$ at an angle where $b_j\ne0$.

In the radial case $\phi=f(u)$, $u=|z|^2$, the hole is a disk by
uniqueness and rotational invariance. Its exact center satisfies
\[
 m u_*f'(u_*)=q.
\]
The asymptotic mode variance is $m u\rho(u)$, with
$\rho(u)=f'(u)+uf''(u)$. The corresponding Gaussian argument is
\begin{equation}\label{eq:radial-growing-jet-argument}
 \frac{m u f'(u)-q}{\sqrt{m u\rho(u)}}.
\end{equation}
For bounded arguments, Taylor expansion at $u_*$ identifies it with
$2\sqrt{m\rho(u_*)}(\sqrt u-\sqrt{u_*})+o(1)$, so
\Cref{thm:curvature-boundary-jets} proves this radial formula for all sublinear
orders. For $f=u^r+c u^{r+1}+\cdots$, its first raw-center shift is
$c(r+1)r^{-3/2}\sqrt q R^2$, with critical growth $m^{2/(r+2)}$.
This larger radial threshold does not describe arbitrary coordinates.

No fixed finite Taylor jet determines the correction for every
$q=o(m)$. For example, near zero take
\[
 f(u)=u^r+c u^r \mathrm{e}^{-1/\sqrt u},\qquad c>0,
\]
smoothly extended at zero and interpolated globally as above.
It has the same Taylor jet as $u^r$ at every finite order, but at
$u=R^2$ its relative contribution to $uf'(u)$ is
$c \mathrm{e}^{-1/R}(1+1/(2rR))$. For $q=\lfloor m/\log m\rfloor$,
$\sqrt q \mathrm{e}^{-1/R}/R\to\infty$. The raw observation with $\tau=0$
is outside the true center by many transition widths, giving limit one.
The exact quadrature-domain correction includes this flat contribution.

